\section{Two complementary uniform estimates}
Write $E_{n,k}(t)$ for the minimum-cycle-length excedance enumerator and put
$P_n=E_{n,k}/t$ at fixed $k$.  We give sufficient conditions for all the
zeros of $P_n$ and $P_{n+1}$ to be simple and negative and to interlace strictly.
In the middle of the root range use
$t=-\exp(-\pi\cot\theta)$, $S=\sin\theta$, and
$z_p=(\cos\theta+ipS)^{-1}$ for odd integers $p$.

\paragraph{Large thresholds.}
For $n\ge10{,}000$ it is sufficient that $k\ge2000\sqrt n$.
The earlier estimate had an additional factor growing logarithmically with
$n/k$. To remove it, let
$D(u)=\exp(\sum_{j\ge k}u^j/j)$ and
$b_j=[u^j](1-u)(D(u)-1)$.  The permutation cycle recurrence gives
$j b_j=-\sum_{i=j-k+1}^{j-1}b_i$ for $j>k$.
Hence on the block $hk<j\le(h+1)k$ we have
$|b_j|\le1/(h!k^2)$ for $h\ge1$.
This gives a summable exponential majorant for the coefficients.

Expand the product $\prod_p D(z_pu)$ by partial fractions, distinguishing
one pole $a$ at a time. Each other occupied pole contributes a numerator
term $b_jz_q^{j-1}u^j/(i(q-a)S)$. For $kS^2\le1/4$ we sum the linear terms in the logarithm
of the product with their signs. Indeed, for any sequence
$f_q=\int e^{-iqv}\,d\mu(v)$ with convergent absolute sum after division by
$q-a$, we have
\[
\left|\sum_{q\ne a}\frac{f_q}{q-a}\right|\le
\frac{\pi}{2}\|\mu\|_{\rm TV}.
\]
This follows by Abel summation of the odd lattice's paired sine series.
The gamma integral supplies the required Fourier measures for both powers
of $z_q$ and their angular derivatives. The higher terms of the logarithm
are bounded absolutely by the square-summable lattice kernel. Together with
the block estimate this controls all occupied-pole counts in this region.
For $kS^2\ge1/4$, reserve powers at the occupied poles and sum the resulting
absolutely convergent majorants. The two estimates give both insertion
signs when $n\ge3k-1$, $nS^2\ge6/5$ and $kS\ge2000$. The exterior-root
theorem supplies the remaining ranks, and the two-cycle theorem covers
$n\le3k-2$.
\href{https://github.com/ProgramComputer/open-math-research/tree/main/topics/long-cycle-excedances/uniform-hilbert}{Full inequalities, exact checks and independent review.}

\paragraph{Shorter thresholds.}
For $n\ge10^6$ and $k\ge400$ it is sufficient that
\[
0.89n/k>47+(5/2)\log n .
\]
Write $w_p=s+ip\pi$, $s=\pi\cot\theta$, $r=|w_1|$, and
$L(w)=\log((1+e^{w-s})/(1+e^{-s}))$ with $L(0)=0$.
The numerator in the coefficient formula is $T_s=e^G$, where $G$ is the
Taylor polynomial of $L$ through degree $k-1$. Compare $G(w)$ on
$r\le|w|\le(1+1/k)r$ with $L(\widehat w)$ at
$\widehat w=(1-1/k)rw/|w|$.  The sum of the two errors at $w$ and $w_1$
is below $34$ in absolute value. For $s\le4\sqrt k$, the absolute Taylor
coefficients of degree $j\ge2$, multiplied by $jr^j$, are at most
$2+9/(4S\sqrt{j})$. Integrating the radial displacement cancels the
harmonic factor. For $s>4\sqrt k$ expand $L$ in exponential frequencies,
bound the first $5k/(4s)$ frequencies directly, and sum a Poisson lower
tail for the remaining ones. The contracted Taylor tail is controlled by
the complementary Poisson upper tail and a geometric series.
Since $|1-\exp(-(s+i\pi)/k)|\ge2/k$, it follows that
$|T_s(w)/T_s(w_1)|<13ke^{34}$.

Choose a common contour with radius between $r(1+0.9/k)$ and $r(1+1/k)$
and at distance at least $\min(r/(40k),1/r)$ from every pole. In the
normalization of the phase-crossing proof the two coefficient errors are at
most $1300n^2\exp(34-0.89n/k)<S/100$. In the normalized difference $Y_{n+1}-(\cos\theta)Y_n$,
the common residue acquires the multiplier $z_p-\cos\theta$, of modulus at most $|p|S$.
The neighboring-residue estimate and exterior-root join then give all the
required next-row signs. These two threshold estimates already cover every
$3\le k\le n$ for $n\ge250000^2$, using the earlier fixed small-threshold
bound for $k<400$.
\href{https://github.com/ProgramComputer/open-math-research/blob/main/topics/long-cycle-excedances/radial-taylor.md}{Proof, independent review and exact constants.}

\subsection{A uniform cutoff of 4,062,720}

\textbf{Theorem.}
For every $n\ge4{,}062{,}720$ and $3\le k\le n$, the two reduced
excedance rows have only simple negative roots and strictly interlace.

\emph{Proof.}
Write $x=n/k$. For $x<4$, use the three-cycle theorem. In the range
$4\le x\le125$, bound the middle angular region $nS^2\ge6/5$ using
the exact pole expansion. For $x\in[\ell,X]$ and $k\ge1000$,
each additional pole $q$ retains at least $k-1$ powers, with
\[
 |z_q|^{k-1}\le w_q=
 \left(1+\frac{999}{160000X}(q^2-1)\right)^{-96}.
\]
For small $X$, marking the two-pole contribution separately controls the
derivative by the larger absolute pole index. For larger $X$, take $M=n$
or $M=n+1-k$ and split according to the total extra numerator degree $T$.
At $T\le M/2$ the distinguished pole retains at least $M/2$ powers.
At $T>M/2$, exponential weighting of the numerator coefficients supplies
$\exp[-(\ell-1)/8]$. In these two sums the available first absolute
moments of $b_j$ are bounded by $6$ and $11$, respectively, so
differentiating an additional pole still gives a convergent positive
majorant. Absolute summation over all subsequent pole choices reduces
the angular estimates to explicit lattice sums. Rational upper bounds on
the tails and the finite sums are sufficient.

More precisely, put $B_j=\sum_p z_p^j$ and
$F_j=r^jE_{j,k}(t)/(j!(1-t)^j)$. At an old zero $t$ the derivative
term $t(1-t)P_n'(t)$ and the insertion lag have the common factor
$n!(1-t)^{n+1}/(t r^{n+1})$ multiplying $SF'_n$ and
$B_k F_{n+1-k}$, respectively. Here primes on $F_n,B_n$ denote angular
derivatives. The quantities $k|F_n-d_nB_n|$,
$k|F'_n-d_nB'_n|/n$, and $k|B_k F_{n+1-k}|/(nS)$ have respective
bounds $V,\Delta,I$ in the
\href{https://github.com/ProgramComputer/open-math-research/tree/main/topics/long-cycle-excedances/adaptive-poles}{proof and exact certificate}.
The single-pole value and derivative tails cost at most $1/50$ and $1/16$.
In each stated parameter interval the certificate proves
\[
(V+1/50)^2+(\Delta+I+1/16+1/100)^2<1.
\]
The principal amplitude is $2kd_n>1$, where $d_n=[u^n]D(u)$ in the
cycle series above. At every old zero the cosine error bound therefore
forces the sine to dominate both the derivative error and the lag.
The old integer-phase samples alternate, and the exterior theorem and
reciprocity give all $n-2$ old roots and the correct next-row signs.

On an aspect interval with upper endpoint $X$ and certified angular
threshold $A_0$ for $kS$, it suffices that
$n\ge\max(10000,1000X,\lceil5(XA_0)^2/6\rceil)$. The complete
18-interval table covers $4\le x\le96$ once $n\ge4{,}062{,}720$,
through $x=100$ at $n\ge4{,}408{,}334$, and through $x=125$ at
$n\ge7{,}500{,}000$. For $400\le k$ the wider-circle criterion covers
the complementary larger $x$ in each of these ranges up to $250000^2$.
Indeed $0.89x>47+2.5\log n$ at the respective joining values
$x=96,100,125$ at their upper row endpoints, using
$\log(4{,}408{,}334)<15.3$, $\log(7{,}500{,}000)<16$, and
$\log(250000^2)<25$. Above $250000^2$ the two complementary uniform
estimates already apply.

For $3\le k<400$ a separate short-circle argument covers every $n\ge10^6$.
For $k\le10$ this follows from the earlier criterion.
For $11\le k<400$ the maximum principle for the reversed Taylor
polynomial bounds $|G|$ on radius $r(1+1/(5k))$ by
$e^{1/5}(6+10\log k)<82$. A contour beyond $r(1+0.15/k)$ gives both
row errors below $S/100$. If $s\le9\sqrt k$ only two poles enter.
Otherwise the gamma-density derivative bound is below $12/(25\pi)$
on the required pole segments, using its Poisson lower tail for
$11\le k\le20$. The same phase comparison and exterior join apply.
This completes the proof.
\href{https://github.com/ProgramComputer/open-math-research/blob/main/topics/long-cycle-excedances/small-threshold-circle.md}{Small-threshold estimates and checks.}

In particular any counterexample must satisfy $k\ge4$ and
$4k\le n<4{,}062{,}720$. The interval estimates exclude additional
pairs inside this region. They do not certify the full finite remainder.
